\section{线性空间*}
在前面的讨论中，向量的运算可以转化为坐标的运算。例如，两个平面向量相加，就是将它们的横、纵坐标分别相加；向量与实数相乘，就是将各个坐标同乘这个实数。这样的运算也可以用于三个、四个乃至任意有限多个数所组成的数组，而不必先为它们找一个几何图形。

函数也有类似的运算。设\(f , g\)都是定义在\(\mathbb{R}\)上的实值函数，规定
\begin{equation}
    \left( f + g \right) \left( x \right) = f \left( x \right) + g \left( x \right) , \quad \left( k f \right) \left( x \right) = k f \left( x \right) .
\end{equation}
容易知道，函数的加法与数乘满足和几何向量相同的运算律。进一步地，一个二次多项式可以写成\(a + b x + c x^{2}\)，正如一个空间向量可以写成\(a \vec{\imath} + b \vec{\jmath} + c \vec{k}\)：都是选取几个固定的对象，将它们分别乘以一个数再相加。

数组、函数与几何向量是不同的对象，但只要研究中用到的是共同的运算规律，就可以用同样的方法处理。将这些规律提炼出来，就得到线性空间的概念。

\subsection{线性空间的定义}
\begin{definition}\label{def:linspace}
设\(V\)是一个集合，在\(V\)上规定了加法与数乘：对任意\(a , b \in V\)，有\(a + b \in V\)；对任意实数\(k\)和\(a \in V\)，有\(k a \in V\)。如果对任意\(a , b , c \in V\)和实数\(k , s\)，下列运算律成立：

（1）\(a + b = b + a\)；

（2）\(\left( a + b \right) + c = a + \left( b + c \right)\)；

（3）存在零向量\(\vec{0}\)，使\(a + \vec{0} = a\)；

（4）每个\(a\)都有负向量\(- a\)，使\(a + \left( - a \right) = \vec{0}\)；

（5）\(1 a = a\)；

（6）\(k \left( s a \right) = \left( k s \right) a\)；

（7）\(k \left( a + b \right) = k a + k b\)；

（8）\(\left( k + s \right) a = k a + s a\)，

则称\(V\)为一个实线性空间，也称实向量空间，\(V\)中的元素统称为向量。
\end{definition}

这里的“向量”有了更一般的含义，它可以是几何向量，也可以是数组、函数等对象。线性空间只规定加法与数乘及其运算律，长度与夹角则需要另行规定。本节和下一节的数乘都以实数为标量。

全体\(n\)元实数组组成的集合记为\(\mathbb{R}^{n}\)，按分量规定加法与数乘：
\begin{equation}
    \left( x_{1} , \ldots , x_{n} \right) + \left( y_{1} , \ldots , y_{n} \right) = \left( x_{1} + y_{1} , \ldots , x_{n} + y_{n} \right) ,
\end{equation}
\begin{equation}
    k \left( x_{1} , \ldots , x_{n} \right) = \left( k x_{1} , \ldots , k x_{n} \right) .
\end{equation}
由实数的运算律，容易验证它是一个线性空间。平面向量与空间向量的坐标运算，分别就是\(\mathbb{R}^{2}\)和\(\mathbb{R}^{3}\)中的运算。

定义在\(\mathbb{R}\)上的全体实值函数，在本节开头规定的运算下，也构成线性空间。其中零向量是恒等于零的函数，\(f\)的负向量是函数\(- f\)。类似地，全体数列也构成线性空间，其运算按每一项进行。

{\kaiti{}思考\quad 只用上述运算律证明：零向量与负向量都是唯一的，并且\(0 a = \vec{0}\)，\(\left( - 1 \right) a = - a\)。}

\subsection{子空间}
平面上过原点的一条直线，在向量的加法与数乘下，自身也构成线性空间；空间中过原点的平面也有同样的性质。这些较小的线性空间都包含在原来的线性空间中。

\begin{definition}\label{def:subspace}
设\(U\)是线性空间\(V\)的一个子集。如果\(U\)在\(V\)的加法与数乘下也构成线性空间，则称\(U\)为\(V\)的子空间。
\end{definition}

判断子空间，只需验证它含零向量，并且对加法、数乘封闭。其余运算律由\(V\)中的运算律直接得到，负向量则由\(- a = \left( - 1 \right) a\)得到。因此，判据也可以写成：\(U\)非空，且对任意\(a , b \in U\)与实数\(k , s\)，都有\(k a + s b \in U\)。

例如，在\(\mathbb{R}^{3}\)中，满足\(x + y - z = 0\)的数组组成一个子空间，几何上就是一个过原点的平面。把方程右边改成\(1\)，所得平面不含零向量，也就不再是子空间。

一般地，形如
\begin{equation}
    \left\{ \begin{matrix}
        a_{11} x_{1} + \cdots + a_{1 n} x_{n} = b_{1} , \\
        a_{21} x_{1} + \cdots + a_{2 n} x_{n} = b_{2} , \\
        \cdots \cdots \cdots \\
        a_{m 1} x_{1} + \cdots + a_{m n} x_{n} = b_{m}
    \end{matrix} \right.
\end{equation}
的方程组称为{\heiti{}线性方程组}。其中\(x_{1} , \ldots , x_{n}\)是未知量，\(a_{i j}\)为系数，\(b_{i}\)为常数项。常数项全为零时，称为{\heiti{}齐次线性方程组}。

齐次线性方程组的解集是\(\mathbb{R}^{n}\)的一个子空间。事实上，若\(u , v\)都是解，则对每个方程都有
\begin{equation}
    \sum_{j = 1}^{n} a_{i j} \left( u_{j} + v_{j} \right)
    = \sum_{j = 1}^{n} a_{i j} u_{j} + \sum_{j = 1}^{n} a_{i j} v_{j} = 0 .
\end{equation}
所以\(u + v\)仍是解，同理\(k u\)也是解，而零向量显然是解。

函数空间中的例子也很多。定义在\(\mathbb{R}\)上的全体奇函数、全体偶函数分别构成子空间；次数不超过\(2\)的实系数多项式连同零多项式也构成一个子空间，记为\(P_{2}\)。这些集合对加法与数乘的封闭性都容易验证。

{\kaiti{}思考\quad 两个子空间的交集还是子空间吗？它们的并集呢？}

\subsection{线性组合与线性无关}
设\(a_{1} , \ldots , a_{m}\)是线性空间中的一组向量。形如
\begin{equation}
    k_{1} a_{1} + k_{2} a_{2} + \cdots + k_{m} a_{m}
\end{equation}
的向量称为这组向量的一个{\heiti{}线性组合}。若\(a\)可以写成上述形式，就称\(a\)可以由\(a_{1} , \ldots , a_{m}\)线性表示。

记\(\operatorname{span}\left( a_{1} , \ldots , a_{m} \right)\)为这组向量的一切线性组合组成的集合，称为由它们{\heiti{}张成的子空间}。两个线性组合相加，或将一个线性组合乘以实数，所得结果仍是线性组合，因此子空间判据直接成立。

在\(\mathbb{R}^{3}\)中，取\(a_{1} = \left( 1 , 0 , 1 \right)\)，\(a_{2} = \left( 0 , 1 , 1 \right)\)。要将\(\left( x , y , z \right)\)表示为\(k_{1} a_{1} + k_{2} a_{2}\)，就是求解
\begin{equation}
    k_{1} = x , \quad k_{2} = y , \quad k_{1} + k_{2} = z .
\end{equation}
所以能被表示的向量恰好满足\(x + y - z = 0\)，也就是前面的平面子空间。可见，求线性表示与解线性方程组是同一个计算问题。

有了表示，进一步的问题就是表示是否唯一。如果同一个向量有两种表示，
\begin{equation}
    k_{1} a_{1} + \cdots + k_{m} a_{m} = s_{1} a_{1} + \cdots + s_{m} a_{m} ,
\end{equation}
相减可得
\begin{equation}
    \left( k_{1} - s_{1} \right) a_{1} + \cdots + \left( k_{m} - s_{m} \right) a_{m} = \vec{0} .
\end{equation}
因此，只需研究零向量的表示：如果零向量只能由系数全为零的组合表示，任意向量的表示就都是唯一的。

\begin{definition}\label{def:linindep}
设\(a_{1} , \ldots , a_{m}\)是线性空间中的一组向量。如果
\begin{equation}
    k_{1} a_{1} + \cdots + k_{m} a_{m} = \vec{0}
\end{equation}
只能在\(k_{1} = \cdots = k_{m} = 0\)时成立，则称这组向量线性无关；否则称它们线性相关。
\end{definition}

线性相关意味着存在一组不全为零的系数使上述等式成立。取其中\(k_{j} \ne 0\)的一项，移项并除以\(k_{j}\)，即可将\(a_{j}\)表示为其余向量的线性组合。反过来，若其中一个向量可由其余向量线性表示，移项就得到一个系数不全为零的零组合。因此，一组向量线性相关，当且仅当其中某个向量可由其余向量线性表示。

这也说明，线性相关的向量组可以删去其中某个向量，而不改变张成的子空间。例如，向上面的向量组中加入\(a_{3} = \left( 1 , 1 , 2 \right)\)，由于\(a_{3} = a_{1} + a_{2}\)，所能表示的向量没有增加，表示却不再唯一：
\begin{equation}
    k_{1} a_{1} + k_{2} a_{2}
    = \left( k_{1} - t \right) a_{1} + \left( k_{2} - t \right) a_{2} + t a_{3} .
\end{equation}

在几何向量中，两个向量线性相关就是共线，三个向量线性相关就是共面。在坐标中，判断线性无关就是判断相应的齐次线性方程组是否只有零解。两种语言描述的是同一个关系。

\begin{example}
（1）判断向量组\(\left( 1 , 0 , 1 \right)\)，\(\left( 0 , 1 , 1 \right)\)，\(\left( 1 , 1 , 0 \right)\)是否线性无关。

（2）在函数空间中，证明\(\sin x , \cos x\)线性无关，而\(1 , \sin^{2} x , \cos^{2} x\)线性相关。

（3）设\(a , b , c\)线性无关，证明\(a + b , b + c , c + a\)线性无关。逆命题是否成立？
\end{example}

\subsection{基与坐标}
为了用一组向量表示整个空间，需要它们张成这个空间；为了让表示唯一，又需要它们线性无关。将这两个要求放在一起，就得到基的概念。

\begin{definition}\label{def:basis}
设\(e_{1} , \ldots , e_{n}\)是线性空间\(V\)中的一组向量。如果它们线性无关，并且张成\(V\)，则称它们为\(V\)的一组基。
\end{definition}

由前面的讨论，\(e_{1} , \ldots , e_{n}\)是一组基，当且仅当\(V\)中的每个向量\(a\)都可以唯一地写成
\begin{equation}
    a = x_{1} e_{1} + x_{2} e_{2} + \cdots + x_{n} e_{n} .
\end{equation}
这里唯一确定的数组\(\left( x_{1} , \ldots , x_{n} \right)\)称为\(a\)在这组基下的坐标。

前面选取互相垂直的单位向量建立坐标，是这里的一种特殊情形。建立坐标的基本要求是表示的存在性与唯一性，因此也可以用其他的基。例如，在\(\mathbb{R}^{2}\)中取
\(e_{1} = \left( 1 , 1 \right)\)，\(e_{2} = \left( 1 , 2 \right)\)。由
\begin{equation}
    \left( x , y \right) = s e_{1} + t e_{2}
\end{equation}
得到\(s + t = x\)，\(s + 2 t = y\)，故\(s = 2 x - y\)，\(t = y - x\)。任意向量都有唯一的表示，所以它们是一组基。\(\left( 3 , 4 \right) = 2 e_{1} + e_{2}\)在这组基下的坐标是\(\left( 2 , 1 \right)\)，在通常的正交单位基下则是\(\left( 3 , 4 \right)\)。同一个向量，换一组基，坐标一般也随之改变。

在\(P_{2}\)中，\(1 , x , x^{2}\)是一组基，多项式\(a + b x + c x^{2}\)的坐标是\(\left( a , b , c \right)\)。取定一组基，就将多项式的加法、数乘转化为系数的运算。

一般地，如果\(a , b\)的坐标分别为\(\left( x_{1} , \ldots , x_{n} \right)\)与\(\left( y_{1} , \ldots , y_{n} \right)\)，则
\begin{equation}
    k a + s b = \sum_{i = 1}^{n} \left( k x_{i} + s y_{i} \right) e_{i} .
\end{equation}
因此它的坐标就是\(\left( k x_{1} + s y_{1} , \ldots , k x_{n} + s y_{n} \right)\)。这说明，取定一组基之后，向量的加法与数乘都化为\(\mathbb{R}^{n}\)中的运算。

{\kaiti{}思考\quad 在上面的基\(e_{1} , e_{2}\)下，若两个向量的坐标分别为\(\left( s_{1} , t_{1} \right)\)与\(\left( s_{2} , t_{2} \right)\)，它们的数量积还是\(s_{1} s_{2} + t_{1} t_{2}\)吗？请写出正确的式子。}

\subsection{维数}
同一个空间可以有不同的基，那么一组基所含向量的个数会不会改变？平面的基总是含两个向量，空间的基总是含三个向量。这种不随基改变的性质，在一般线性空间中也成立。

\begin{theorem}\label{thm:dim}
设\(e_{1} , \ldots , e_{n}\)与\(f_{1} , \ldots , f_{m}\)是同一个线性空间的两组基，则\(n = m\)。
\end{theorem}

{\heiti{}证明思路：}
\nopagebreak[4]

线性无关的向量组没有可以删去的向量，而张成空间的向量组已经足以表示所有向量。因此，如果一组线性无关的向量都能由另一组向量表示，前者的个数就不应超过后者。证明这个事实，可以逐个将线性无关的向量换入原来的向量组，并保持所张成的空间不变。

{\heiti{}证明：}
\nopagebreak[4]

设\(u_{1} , \ldots , u_{p}\)线性无关，且都能由\(w_{1} , \ldots , w_{q}\)表示。记\(W = \operatorname{span}\left( w_{1} , \ldots , w_{q} \right)\)。先将\(u_{1}\)表示为\(w\)组的线性组合，由于\(u_{1} \ne \vec{0}\)，至少有一个系数不为零。不妨设\(w_{1}\)的系数不为零，则可以反过来将\(w_{1}\)表示为\(u_{1} , w_{2} , \ldots , w_{q}\)的线性组合。因此，用\(u_{1}\)替换\(w_{1}\)后，张成的空间仍是\(W\)。

假设已经换入\(u_{1} , \ldots , u_{k - 1}\)，现在将\(u_{k}\)用当前向量组表示。如果剩下的\(w\)的系数全为零，就将\(u_{k}\)表示为前面几个\(u\)的线性组合，与线性无关矛盾。因此仍有某个\(w\)的系数不为零，可以将它替换为\(u_{k}\)，并保持张成的空间不变。

每次替换，\(u\)多一个而\(w\)少一个。如果\(p > q\)，经过\(q\)次替换后，\(u_{1} , \ldots , u_{q}\)已经张成\(W\)，于是\(u_{q + 1}\)可由它们表示，仍与线性无关矛盾。故\(p \le q\)。

现在将这个事实用于两组基。因为\(e\)组线性无关，且每个\(e_{i}\)都可由\(f\)组表示，所以\(n \le m\)；交换两组的角色，又得\(m \le n\)，故\(n = m\)。

能由有限个向量张成的线性空间称为{\heiti{}有限维线性空间}。从一组有限的张成向量中逐次删去可以由其余向量表示的向量，就得到一组基。这组基中向量的个数称为空间的{\heiti{}维数}，记为\(\dim V\)。上面的定理保证维数不依赖于基的选择。

例如，\(\mathbb{R}^{n}\)的标准基
\begin{equation}
    e_{1} = \left( 1 , 0 , \ldots , 0 \right) , \quad
    e_{2} = \left( 0 , 1 , \ldots , 0 \right) , \quad \ldots , \quad
    e_{n} = \left( 0 , \ldots , 0 , 1 \right)
\end{equation}
含有\(n\)个向量，所以\(\dim \mathbb{R}^{n} = n\)。\(P_{2}\)的基\(1 , x , x^{2}\)含三个向量，所以它也是三维空间。前面平面\(x + y - z = 0\)中的\(a_{1} , a_{2}\)是一组基，因此这个子空间的维数是\(2\)。维数所计数的是一组基中的向量，也就是描述空间中任意向量所需的坐标个数。

在\(n\)维空间中，线性无关的向量组至多含\(n\)个向量。若已有一组线性无关的向量而它们尚未张成整个空间，就选取一个不能由它们表示的向量加入，所得向量组仍线性无关。这样继续添加，至多达到\(n\)个向量，就得到一组基。因此，任意线性无关的向量组都可以补成一组基，特别地，\(n\)个线性无关的向量已经是一组基。

对子空间也可以这样选取向量，至多经过\(n\)次就得到它的一组基。因此，有限维空间的子空间仍然有限维，并且维数不超过原空间。

也有不能由有限个向量张成的线性空间，称为{\heiti{}无限维线性空间}。例如，全体实系数多项式组成的空间就是无限维的：任取有限个多项式，总能找到一个非负整数\(d\)，使它们的次数都不超过\(d\)，则它们的任何线性组合都不能得到\(x^{d + 1}\)。所以任意有限组多项式都不能张成整个多项式空间。

{\kaiti{}思考\quad 定义在\(\mathbb{R}\)上的全体实值函数组成的空间，是有限维的吗？}

\begin{exercise}
（1）判断下列\(\mathbb{R}^{3}\)的子集是否为子空间：\\
（a）\(\left\{ \left( x , y , z \right) \mid x + y + z = 0 \right\}\)；\\
（b）\(\left\{ \left( x , y , z \right) \mid x + y + z = 1 \right\}\)；\\
（c）\(\left\{ \left( x , y , z \right) \mid x y = 0 \right\}\)；\\
（d）\(\left\{ \left( x , y , z \right) \mid x = 2 y , z = 3 y \right\}\)。\\
对其中的子空间，求一组基及其维数。

（2）设\(a_{1} = \left( 1 , 0 , 1 \right)\)，\(a_{2} = \left( 0 , 1 , 1 \right)\)，\(a_{3} = \left( 1 , 1 , t \right)\)。讨论\(t\)取何值时，这组向量线性无关。

（3）在\(\mathbb{R}^{2}\)中，验证\(e_{1} = \left( 1 , 2 \right)\)，\(e_{2} = \left( 1 , 3 \right)\)构成一组基，并求\(\left( 4 , 7 \right)\)在这组基下的坐标。

（4）判断\(\left( 17 , - 4 , 2 \right)\)和\(\left( 17 , - 4 , 5 \right)\)能否由\(\left( 2 , 1 , - 3 \right)\)，\(\left( 1 , - 2 , 4 \right)\)线性表示；如果能，求出表示式。

（5）求齐次线性方程组
\begin{equation}
    \left\{ \begin{matrix}x + y + z = 0 , \\ x - y + z = 0\end{matrix} \right.
\end{equation}
的解集的一组基及其维数。

（6）在\(P_{2}\)中，判断\(1 , x + 1 , \left( x + 1 \right)^{2}\)是否构成一组基。若构成，求\(x^{2}\)在这组基下的坐标。

（7）证明：每个定义在\(\mathbb{R}\)上的实值函数都可以唯一地表示为一个奇函数与一个偶函数之和。

（8）在正实数集上规定\(u \oplus v = u v\)，\(k \odot u = u^{k}\)。验证它在这两种运算下是实线性空间，并求一组基及其维数。
\end{exercise}
\newpage
